\documentclass[11pt]{article}
\usepackage[a4paper,margin=1in]{geometry}
\usepackage{amsmath,amssymb,amsthm}
\usepackage{booktabs}
\usepackage{array}
\usepackage{hyperref}

\newtheorem{definition}{Definition}
\newtheorem{theorem}{Theorem}

\newcommand{\Hh}{\mathbb H}
\newcommand{\GammaS}{\Gamma}
\newcommand{\XiGamma}{\Xi_{\Gamma}}
\newcommand{\ZG}{Z_{\Gamma}}
\newcommand{\CRE}{\mathrm{CRE}}
\newcommand{\CRCFT}{\mathrm{CRCFT}}

\title{Step 185: Selberg Trace Formula / Maass Carrier Pivot}
\author{RH Membrane Adequacy Track}
\date{}

\begin{document}
\maketitle

\section*{Verdict}

\[
\boxed{\mathrm{selberg\_verdict}=V_{\mathrm{selberg\_closes\_with\_proved\_RH}}.}
\]

The Selberg/Maass carrier is RH-analogous but not classically-RH-equivalent
for Riemann's zeta.  On its native carrier, the framework residual closes:
\[
\boxed{\XiGamma=0}
\]
after the full Selberg spectral/geometric ledger is included.  This closure
does not transfer to Riemann RH.

\section*{Inherited Record Audit}

\paragraph{Step 90.}
The source-family audit says that trace-formula and automorphic sources may
provide genuine spectral/geometric positivity only if the completed carrier is
chosen to be automorphic/trace-native and if Plancherel or trace-formula data
provide the full lower-frame record.  Otherwise such data are public trace
shadows.

\paragraph{Step 92.}
The Hecke carrier sketch records Selberg/trace-formula sources as plausible
only if carrier-native and accompanied by explicit bridge records.  A trace
formula cannot be inserted as a small public repair term into a different
zeta carrier.

\paragraph{Consequence for this step.}
Unlike previous pivots, this step chooses the trace-native carrier itself.
Thus the Selberg trace formula is not used as a public shadow for Riemann's
zeta; it is used only to close the native Selberg residual.

\section*{Carrier Declaration}

Let \(\Gamma=\mathrm{SL}_2(\mathbb Z)\), with the usual
\(\mathrm{PSL}_2(\mathbb Z)\) quotient understood for the action on
\(\Hh\).  Define
\[
H_\Gamma=L^2(\Gamma\backslash\Hh,d\mu),
\qquad
d\mu=\frac{dx\,dy}{y^2}.
\]
The native operator is the positive hyperbolic Laplacian
\[
\Delta=-y^2\left(\partial_x^2+\partial_y^2\right),
\]
with its standard self-adjoint realization on \(H_\Gamma\).

Its spectral decomposition contains:
\begin{itemize}
\item discrete Maass cusp-form eigenvalues
\[
\lambda_j=\frac14+r_j^2;
\]
\item residual/discrete terms;
\item continuous Eisenstein spectrum, with scattering/parabolic corrections.
\end{itemize}

\section*{Selberg Zeta}

The Selberg zeta function is
\[
\ZG(s)
=
\prod_{P\ {\rm primitive}}
\prod_{k\ge0}
\left(1-N(P)^{-s-k}\right),
\]
where \(P\) ranges over primitive closed geodesics and \(N(P)\) is the
geodesic norm.  Spectral zeros associated to the Laplacian satisfy
\[
s(1-s)=\lambda_j.
\]
For \(\lambda_j=1/4+r_j^2\) this gives
\[
s=\frac12\pm ir_j.
\]
Residual, exceptional, parabolic, and continuous-spectrum terms are native
Selberg ledger terms, not omissions.

\section*{Parent Residual}

\begin{definition}[Selberg adequacy residual]
\(\XiGamma\) is the Schur/trace-formula defect between:
\begin{itemize}
\item native spectral probes: the discrete Maass spectrum, residual terms, and
Eisenstein/scattering spectral measure of \(\Delta\);
\item native geometric probes: identity, hyperbolic closed-geodesic terms,
parabolic terms, and scattering corrections.
\end{itemize}
The audit energy is the Selberg transform/Plancherel trace-test pairing.
\end{definition}

Equivalently, for admissible test functions \(h\), write the Selberg trace
formula schematically as
\[
\mathcal S_\Gamma(h)=\mathcal G_\Gamma(h).
\]
Then
\[
\XiGamma(h):=\mathcal S_\Gamma(h)-\mathcal G_\Gamma(h).
\]

\begin{theorem}[Native Selberg closure]
On the Selberg/Maass carrier for \(\Gamma=\mathrm{SL}_2(\mathbb Z)\), after
all discrete, continuous, residual, identity, hyperbolic, parabolic, and
scattering terms are included,
\[
\XiGamma=0.
\]
\end{theorem}

\begin{proof}
This is precisely the Selberg trace formula on the native automorphic
carrier.  The spectral side is the spectral decomposition of the self-adjoint
Laplacian on \(L^2(\Gamma\backslash\Hh)\).  The geometric side is the closed
geodesic/orbital expansion plus identity, parabolic, scattering, and residual
terms.  Since both sides are the same distribution on the admissible test
class, the Schur/trace residual vanishes.

The Selberg zeta RH analog is then native: the nontrivial spectral zeros tied
to self-adjoint Laplacian eigenvalues lie on the line \(s=1/2+ir\) through
\(s(1-s)=\lambda\), with the non-cuspidal and residual terms accounted in the
same ledger.
\end{proof}

\section*{CRE and CRCFT Audit}

\begin{center}
\begin{tabular}{p{0.28\linewidth}p{0.18\linewidth}p{0.44\linewidth}}
\toprule
Audit & Result & Reason\\
\midrule
CRE status & not \(\CRE\)
& Closing \(\XiGamma\) for \(\ZG\) does not imply Riemann RH and is not
logically equivalent to it.\\
\(\CRCFT\) applicability & does not apply
& \(\CRCFT\) is a foreclosure taxonomy for CRE carriers; this carrier closes
natively instead.\\
Framework closure & closed native
& Full Selberg trace formula supplies the exact spectral/geometric ledger.\\
Transfer to Riemann RH & not claimed
& A separate bridge from Selberg/Maass data to Riemann's zeta carrier is not
supplied.\\
\bottomrule
\end{tabular}
\end{center}

\section*{Implication}

The Selberg/Maass pivot tests the boundary of \(\CRCFT\).  The result is that
the CRE hypothesis is load-bearing.  A proved RH-analogous carrier with its
own native trace formula closes under the framework, while the previous CRE
carriers foreclosed because their closure would decide Riemann RH or require
missing carrier bridges.

\section*{Classical Theorems Cited}

\begin{itemize}
\item A. Selberg, 1956: Selberg trace formula and Selberg zeta theory.
\item D. Hejhal, \emph{The Selberg Trace Formula for PSL(2,R)}, Vols. I--II:
standard trace formula, zeta, scattering, and continuous-spectrum treatment.
\item H. Iwaniec, \emph{Spectral Methods of Automorphic Forms}: Maass forms,
Eisenstein spectrum, and automorphic spectral decomposition for
\(\mathrm{SL}_2(\mathbb Z)\).
\end{itemize}

\section*{Nonclaim Boundary}

This step does not prove Riemann RH, does not transfer Selberg's RH analog to
Riemann's zeta, and does not close any prior Riemann-zeta residual.  It closes
only the native Selberg residual \(\XiGamma\) on the Selberg/Maass carrier
with its full trace-formula ledger.

\end{document}
